\section{Bab IV: Dinamika Pelaksanaan UUD NRI Tahun 1945}
\subsection{A. Pengantar Konstitusi}
Istilah konstitusi berasal dari bahasa Prancis (\textit{constituer}) dan Inggris (\textit{constitution}) yang berarti "membentuk suatu negara". UUD NRI Tahun 1945 berkedudukan sebagai konstitusi tertulis (\textit{documentary constitution}) dan merupakan sumber hukum tertinggi. Segala peraturan perundang-undangan di Indonesia harus bersumber pada Pancasila dan UUD 1945.

\subsection{B. Pembabakan Dinamika Pelaksanaan UUD}
\begin{description}
    \item[Periode Awal Kemerdekaan (1945--1949):] Pelaksanaan UUD belum optimal karena fokus mempertahankan kemerdekaan. Terjadi perubahan sistem pemerintahan dari Presidensial ke Parlementer melalui Maklumat 14 November 1945, diwarnai pergantian beberapa kabinet.
    \item[Masa Konstitusi RIS (1949--1950):] Bentuk negara berubah menjadi Serikat/Federasi dengan sistem pemerintahan Parlementer.
    \item[Masa UUDS 1950 (1950--1959):] Kembali ke bentuk Negara Kesatuan, namun menggunakan sistem pemerintahan Parlementer yang menyebabkan ketidakstabilan politik (7 kali pergantian kabinet).
    \item[Kembali ke UUD 1945 (Orde Lama, 1959--1966):] Ditegaskan melalui \textbf{Dekrit Presiden 5 Juli 1959}. Pada masa ini berlaku Demokrasi Terpimpin, dan terjadi penyimpangan seperti penetapan Presiden Seumur Hidup.
    \item[Masa Orde Baru (1966--1998):] Bermaksud melaksanakan Pancasila dan UUD 1945 secara murni dan konsekuen, namun dalam praktiknya terjadi kecenderungan kekuasaan terpusat (\textit{executive heavy}) dan maraknya KKN.
\end{description}